\chapter*{ОПРЕДЕЛЕНИЯ, ОБОЗНАЧЕНИЯ И СОКРАЩЕНИЯ}
\addcontentsline{toc}{chapter}{ОПРЕДЕЛЕНИЯ, ОБОЗНАЧЕНИЯ И СОКРАЩЕНИЯ}

\renewcommand{\arraystretch}{1.15}
\begin{longtable}{@{}>{\justifying\parindent0pt\arraybackslash}p{0.28\textwidth}@{\hspace{8mm}}>{\justifying\parindent0pt\arraybackslash}p{0.62\textwidth}@{}}
\termrow{РС}{рекомендательная система, программа, ранжирующая объекты каталога по предполагаемой полезности для пользователя}
\termrow{PCv1, PCv2, PCv3 (v1, v2, v3)}{версии контентной РС с входом в виде одного объекта-затравки, набора лайков, лайков и дизлайков}
\termrow{Лайк, дизлайк}{явная положительная или отрицательная оценка объекта пользователем}
\termrow{Близость $\mathrm{sim}(x,y)$}{мера сходства двух объектов, $\mathrm{sim}\in[0,1]$, $\mathrm{sim}(x,x)=1$}
\termrow{Расстояние $d(x,y)$}{мера различия двух объектов, в работе $d = 1-\mathrm{sim}$}
\termrow{Матрица смежности $G$}{матрица $n\times n$ попарных значений меры, $G_{ij}=\mathrm{sim}(o_i,o_j)$}
\termrow{LCA}{от англ. lowest common ancestor, наиболее глубокий общий предок двух узлов дерева}
\termrow{SMC}{от англ. simple matching coefficient, коэффициент простых совпадений Сокала--Миченера}
\termrow{Таксономия вкуса}{авторская иерархия <<Ликёры $\rightarrow$ L1 $\rightarrow$ L2 $\rightarrow$ L3>>, классифицирующая объекты по вкусу}
\termrow{Группа интересов}{подмножество лайков, попарно связанных цепочкой близостей не ниже порога $\theta$}
\termrow{Затравка}{объект, по которому строится список рекомендаций в PCv1}
\termrow{Функция полезности $U(x)$}{взвешенная степень выполнения фильтров поискового запроса объектом $x$}
\termrow{recall(N)}{доля скрытых отложенных объектов, попавших в первые $N$ позиций выдачи}
\termrow{coverage(N)}{доля групп интересов пользователя, представленных в первых $N$ позициях выдачи}
\termrow{Hit(k)}{доля сценариев, в которых искомый объект занял место не ниже $k$}
\end{longtable}
